% Implementation Section for AetherGuard SCI Paper
% Target: IEEE TDSC/TIFS

\section{Implementation}
\label{sec:implementation}

\subsection{Technology Stack}
\label{subsec:tech-stack}

AetherGuard is implemented using a modern, type-safe technology stack chosen for performance, maintainability, and operational reliability.

\subsubsection{Control Plane}

\begin{itemize}
    \item \textbf{Runtime}: Node.js 18 LTS with TypeScript 5.0
    \item \textbf{Web Framework}: Express.js 4.18 for REST APIs
    \item \textbf{Real-time Communication}: Socket.IO 4.6 for WebSocket gateways
    \item \textbf{Database}: SQLite 3.42 (default) or PostgreSQL 15 (production)
    \item \textbf{Cache/Message Broker}: Redis 7.0
    \item \textbf{Authentication}: JWT with RS256 signing, bcrypt password hashing
\end{itemize}

\subsubsection{Distributed Agents}

\begin{itemize}
    \item \textbf{Runtime}: Node.js 18 LTS with TypeScript 5.0
    \item \textbf{Process Management}: PM2 for production deployment
    \item \textbf{Network Monitoring}: libpcap via node-pcap bindings
    \item \textbf{Firewall Control}: iptables/nftables via child process invocation
    \item \textbf{System Integration}: systemd service files for Linux deployment
\end{itemize}

\subsubsection{AI/ML Components}

\begin{itemize}
    \item \textbf{Local Models}: Ollama 0.1.x with Llama 3 8B and Mistral 7B
    \item \textbf{Cloud Models}: GLM-4 (Zhipu AI), GPT-4 (OpenAI) via REST APIs
    \item \textbf{Threat Classification}: scikit-learn Random Forest (Python bridge)
    \item \textbf{Embeddings}: sentence-transformers for semantic similarity
\end{itemize}

\subsubsection{Frontend Dashboard}

\begin{itemize}
    \item \textbf{Framework}: React 18 with TypeScript
    \item \textbf{Build Tool}: Vite 5.0
    \item \textbf{UI Components}: TailwindCSS 3.4, Headless UI
    \item \textbf{Charts}: Recharts for real-time metrics visualization
    \item \textbf{Real-time Updates}: Socket.IO client
\end{itemize}

\subsection{Key Implementation Details}
\label{subsec:implementation-details}

\subsubsection{Agent Registration Flow}

The agent onboarding process implements cryptographic verification:

\begin{enumerate}
    \item Administrator generates a one-time onboarding code via the dashboard.
    \item Agent installer presents the code to the \texttt{/api/agents/register} endpoint.
    \item Control plane verifies the code and issues a long-term agent ID and API key.
    \item Agent establishes WebSocket connection with mutual TLS authentication.
    \item Agent receives initial configuration including threat intelligence feeds and policy rules.
\end{enumerate}

\subsubsection{Secure Command Execution}

Commands from the control plane follow a defense-in-depth approach:

\begin{enumerate}
    \item Command is signed with HMAC-SHA256 using the agent-specific signing key.
    \item Command is transmitted over the authenticated WebSocket connection.
    \item Agent verifies the HMAC signature before parsing the command.
    \item Command is validated against an allowlist of permitted operations.
    \item Command execution is sandboxed using Linux namespaces where applicable.
    \item Execution result is signed and returned to the control plane.
\end{enumerate}

\subsubsection{Multi-Agent Debate Implementation}

The debate layer is implemented as an asynchronous orchestration workflow:

\begin{verbatim}
async function debate(threat: Threat): Promise<Decision> {
  // Parallel analysis by role-playing agents
  const viewpoints = await Promise.all([
    socrates.analyze(threat),
    plato.analyze(threat),
    aristotle.analyze(threat)
  ]);

  // Conflict detection
  const conflicts = detectConflicts(viewpoints);

  // Iterative debate rounds
  for (let round = 0; round < MAX_ROUNDS; round++) {
    if (conflicts.length === 0) break;

    const arguments = await Promise.all(
      conflicts.map(c => agents[c.agent].presentArgument(c))
    );

    viewpoints = updateViewpoints(viewpoints, arguments);
    conflicts = detectConflicts(viewpoints);
  }

  // Final synthesis
  return aristotle.synthesize(viewpoints);
}
\end{verbatim}

\subsubsection{Privacy-Aware Model Routing}

The routing implementation uses a decision tree based on threat sensitivity:

\begin{verbatim}
function selectModel(threat: Threat, policy: PrivacyPolicy): Model {
  const sensitivity = assessSensitivity(threat);

  if (policy.mode === 'STRICT' || sensitivity === 'HIGH') {
    return localModel;  // Ollama
  }

  if (sensitivity === 'MEDIUM' && policy.allowHybrid) {
    return hybridModel; // Local pre-screen + cloud verification
  }

  if (policy.costBudget > cloudModel.estimateCost(threat)) {
    return cloudModel;  // GLM-4 or GPT-4
  }

  return localModel;
}
\end{verbatim}

Sensitive data redaction uses regular expression patterns:

\begin{verbatim}
const REDACTION_PATTERNS = {
  ipv4: /\b(?:[0-9]{1,3}\.){3}[0-9]{1,3}\b/g,
  email: /\b[A-Za-z0-9._%+-]+@[A-Za-z0-9.-]+\.[A-Z|a-z]{2,}\b/g,
  token: /bearer\s+[a-zA-Z0-9_\-\.]+/gi,
  apiKey: /api[_-]?key["']?\s*[:=]\s*["']?[a-zA-Z0-9]{16,}["']?/gi
};
\end{verbatim}

\subsection{Deployment Architecture}
\label{subsec:deployment}

\subsubsection{Control Plane Deployment}

The control plane is deployed as a containerized application using Docker Compose:

\begin{verbatim}
version: '3.8'
services:
  backend:
    build: ./backend
    ports:
      - "3005:3005"
    environment:
      - NODE_ENV=production
      - JWT_SECRET=${JWT_SECRET}
      - DATABASE_URL=sqlite:./data/aetherguard.db
    volumes:
      - ./data:/app/data

  redis:
    image: redis:7-alpine
    command: redis-server --requirepass ${REDIS_PASSWORD}

  frontend:
    build: ./frontend
    ports:
      - "8080:80"
\end{verbatim}

\subsubsection{Agent Deployment}

Distributed agents are deployed as systemd services on Linux servers:

\begin{verbatim}
[Unit]
Description=AetherGuard Security Agent
After=network.target

[Service]
Type=simple
User=aetherguard
ExecStart=/usr/bin/node /opt/aetherguard/agent.js
Restart=always
RestartSec=10
Environment=AETHERGUARD_SERVER=https://control.example.com
Environment=AETHERGUARD_AGENT_ID=agent-001

[Install]
WantedBy=multi-user.target
\end{verbatim}

\subsubsection{ESP32 Deployment Dongle}

The ESP32-based deployment dongle provides zero-touch agent installation:

\begin{enumerate}
    \item Dongle connects to target server via USB.
    \item Dongle presents as USB mass storage with auto-run script.
    \item Auto-run script detects OS and downloads appropriate agent binary.
    \item Agent registers with control plane using pre-configured onboarding code.
    \item Dongle LED indicates installation status (green=success, red=failure).
\end{enumerate}

\subsection{Code Availability}
\label{subsec:code}

The AetherGuard implementation is open-source under the MIT license. The repository includes:

\begin{itemize}
    \item Control plane backend source code (Node.js/TypeScript)
    \item Distributed agent source code (Node.js/TypeScript)
    \item Frontend dashboard source code (React/TypeScript)
    \item ESP32 firmware for deployment dongles (C++/Arduino)
    \item Docker Compose configurations for deployment
    \item Systemd service files for agent installation
    \item Comprehensive documentation and API specifications
\end{itemize}

The source code is available at: \url{https://github.com/aetherguard/aetherguard}
